% SPDX-License-Identifier: Apache-2.0
% covers: razboj::raster::Raster
\datasheet{Raster}{Razboj's rasteriser: it reads a display list over AXI and writes one pixel a cycle.}

\dsfacts{\code{gpu/razboj/src/raster.rs}}{\code{razboj::raster::Raster}}%
{\code{//docs:razboj}, \code{//docs:noc}}

\subsection*{Function}

\code{Raster} finds its work in memory and draws it there. While
\code{ring} is high it reads a count at \code{CTRL}, until the count
is not zero. It then reads the sixteen words of each instruction of
the display list at \code{DL}, as one read burst of sixteen beats, and
walks the instruction's box one pixel per cycle. It writes the
instruction's colour, its alpha in the top byte, at every pixel inside
the primitive: every pixel of a clear or a rectangle, and every pixel
of a triangle at which no edge function is negative. A run of a row's
pixels goes out as one write burst of up to sixteen beats, a beat a
pixel, and a beat whose pixel is outside the primitive has its strobes
off (issue 987). A shaded
triangle's colour is three planes the host worked out, each stepped
across the box as an edge is.

It draws list after list. When a list is drawn and every write has
been answered, it writes the count back to zero and reads it again.
On the board the count is \code{Doorbell}'s register, which drives
\code{ring} high while the count is not zero, so an idle rasteriser
puts nothing on the link (issue 985). Where nothing drives it,
\code{ring} is tied high and the count is read back to back.

It is an AXI host written as hardware. It issues bursts on
\code{issue}, sends a pixel's beat on \code{wbeat}, and hands
identifiers back on \code{release}. It receives \code{grant},
\code{done} and \code{rdata} from its tracker, and \code{ring}.
\code{idle} goes high when the list is drawn, its zero is answered and
no write is in flight.

\subsection*{Parameters}

\begin{center}\footnotesize
\begin{tabular}{@{}lp{0.7\textwidth}@{}}
\toprule
Parameter & Meaning \\
\midrule
\code{A} & The width of the link's address, in bits. \\
\code{I} & The width of the AXI identifier, in bits. \\
\code{LOGW} & The screen is \code{1 << LOGW} pixels wide. \\
\code{H} & The screen is \code{H} pixels high. \\
\code{BASE} & The byte address of the framebuffer's first word. \\
\code{DL} & The byte address of the display list's first word. \\
\code{CTRL} & The byte address of the count. \\
\bottomrule
\end{tabular}
\end{center}

The tables use \code{Raster<16, 2, 4, 16, 0, 0x400, 0x600>}, the
rasteriser of \code{//gpu/razboj:trace}: a sixteen by sixteen screen at
\code{0x0}, the list at \code{0x400} and the count at \code{0x600}. In
\code{soc} it is \code{Raster<32, 2, 7, 96, 0x8000, 0x2000, 0x4000>}.
On the flagship it is \code{Raster<32, 2, 10, 480, 0x4200\_0000,
0x4280\_0000, 0x3900>}: rows of 1024 pixels, the scanout's stride, of
which 640 are shown, and the count in \code{Doorbell}.

\subsection*{Address map}

\begin{center}\footnotesize
\begin{tabular}{@{}lp{0.6\textwidth}@{}}
\toprule
Address & What the rasteriser does there \\
\midrule
\code{CTRL} & Reads the count while \code{ring} is high, and writes
zero there when a list is drawn. \\
\code{DL + (i << 6) + (k << 2)} & Reads word \code{k}, 0 to 15, of
instruction \code{i}, as one burst. \\
\code{BASE + (((y << LOGW) + x) << 2)} & Writes pixel \code{x},
\code{y}. \\
\bottomrule
\end{tabular}
\end{center}

An instruction's words are in the format \code{gpu/razboj/src/dl.rs}
states.

\begin{center}\footnotesize
\begin{tabular}{@{}lp{0.42\textwidth}p{0.3\textwidth}@{}}
\toprule
Word & Low field & High field \\
\midrule
0 & bits 1 to 0: kind, 0 clear, 1 rectangle, 2 triangle, 3 shaded
triangle & bits 25 to 2: colour \\
1 & bits 9 to 0: \code{x0} & bits 25 to 16: \code{y0} \\
2 & bits 9 to 0: \code{x1} & bits 25 to 16: \code{y1} \\
3 & bits 15 to 0: \code{ax} & bits 31 to 16: \code{ay} \\
4 & bits 15 to 0: \code{bx} & bits 31 to 16: \code{by} \\
5 & bits 15 to 0: \code{cx} & bits 31 to 16: \code{cy} \\
6 to 14 & \multicolumn{2}{l}{a shaded triangle's red, green and blue
planes: value, step right, step down} \\
15 & bits 7 to 0: alpha & \\
\bottomrule
\end{tabular}
\end{center}

\dstables{Raster}

\subsection*{Behaviour}

The rasteriser is two processes. The first takes the link's answers
every cycle, and says whether the pixel under the walk is in the
primitive and whether the rasteriser is idle. The second is the front
end written as the sequence it is: poll the count until the list is
ready, then for each instruction fetch its sixteen words and walk its
box a pixel a turn, and when the list is drawn write the count back to
zero and poll again. The lowering
numbers the front end's waits into a state register the unit does not
declare, \code{at\_wait}, and keeps its loops' counts in \code{for0} to
\code{for3}; all are in the tables above.

\begin{itemize}
\item The front end waits for \code{ring}, then issues a read of the
count at \code{CTRL} once \code{issue} has room, and waits for its
beat. The count's low sixteen bits go into \code{left}. A count of
zero means wait for \code{ring} and read again.
\item For each of \code{left} instructions it waits a cycle, for the
instruction's index to read back, then reads its sixteen words as one
burst of sixteen beats, issued once \code{issue} has room. Each
word's fields are latched as its beat lands.
\item Once the last word has landed the walk is set up over three
cycles: the differences of the vertices, then the six products side
by side, then their sums. The products are computed there, and
nowhere else.
\item A clear's box is the whole screen, from its own type: all
\code{1 << LOGW} columns of \code{H} rows. A rectangle's and a
triangle's box is \code{x0}, \code{y0} to \code{x1}, \code{y1}, and
is not clipped to \code{H}. Vertices are sixteenths of a pixel in
sixteen bits of two's complement.
\item After a cycle for the box to read back, the walk goes over every
row of the box, a cycle at the start of each, and every column of the
row, a pixel a turn.
\item Each edge is kept as its value at this pixel, its value at the
start of this row and its two steps, in 32-bit two's complement. A
step to the right adds the column step; a new row reloads from the row
value plus the row step.
\item A pixel is sampled at its centre, and a vertex is in sixteenths
of a pixel. \code{hit} is high when the kind is not a triangle, or when
all three edge values are positive, or zero on a top or a left edge:
\code{tl0} to \code{tl2} say which edges are, from the signs of
their steps, set at the setup. So of two triangles sharing an edge,
exactly one draws a pixel whose centre lies on it.
\item A pixel with \code{hit} high and no burst under way starts one,
when both \code{issue} and \code{wbeat} have room, which adds one to
\code{issued}. The burst is size 2, \code{INCR}, and its length is the
pixels left in the box's row, the words left before the next 4\,KB
page, or sixteen, whichever is fewest, so no burst crosses a page.
\code{beats} holds the beats it still owes.
\item While a burst is under way, every pixel the walk steps over sends
the next beat when \code{wbeat} has room, whatever the pixel: the
alpha above the 24-bit colour, strobe \code{0xf} where \code{hit} is
high and \code{0x0} where it is not, and \code{last} on the burst's
last beat. A triangle's pixels in a row are one run, so a burst wastes
beats only past the run's end.
\item The walk steps when the pixel needs no write or its beat went
out, so backpressure holds the walk.
\item When the last instruction's box is walked and every write is
answered, the count at \code{CTRL} is written back to zero, and
\code{finished} is set as that write goes out. \code{inflight} is a wire, \code{issued} less \code{answered},
the writes whose response has not come back, and \code{idle} is high
when \code{finished} is set and nothing is in flight.
\item When a write response and a read beat arrive together, the
write's identifier goes back first and the read's waits a cycle.
Grants are taken and dropped.
\end{itemize}

\subsection*{Verification}

\begin{itemize}
\item \code{//gpu/razboj:razboj\_test} renders scenes on the rasteriser through
an \code{AxiHost}, an \code{AxiPer} and \code{Fb}, and compares every
pixel with the model of \code{gpu/razboj/src/model.rs}:
\code{a\_scene\_of\_every\_kind\_agrees\_with\_the\_model},
\code{a\_clear\_says\_only\_its\_colour},
\code{a\_triangle\_is\_the\_same\_whichever\_way\_it\_is\_wound},
\code{a\_triangle\_hanging\_off}\allowbreak\code{\_the\_screen\_is\_clipped},
\code{a\_single\_pixel\_and\_an\_empty\_box}, and
\code{pseudorandom\_scenes\_agree\_with\_the\_model}, twelve scenes of
up to four random rectangles and triangles.
\item \code{an\_instruction\_survives\_the\_format} in the same target
checks the display list's encoder against its decoder.
\item The build simulates the netlist of \code{Raster<16, 2, 4, 16, 0,
0x400, 0x600>} against the trace of \code{//gpu/razboj:trace} under nvc and
Verilator: \code{//docs:razboj\_sim\_raster\_raster\_tb\_test} and
\code{//docs:razboj\_vsim\_raster\_test}.
\item \code{//soc:demo\_test} runs
\code{a\_core\_draws\_a\_solid\_and\_a\_gpu\_fills\_it}. Vreteno writes
a display list on the lattice, the rasteriser draws it, and every pixel
must match the model's picture of that list.
\item It goes through Vivado out of context with the board's
parameters, \code{//gpu/razboj:razboj\_synth} and
\code{//gpu/razboj:razboj\_pnr}, and meets the board bus's 100~MHz
since its setup was split into three cycles (issue 1034); the size and
the slack of the last run are in \code{//docs:razboj}'s section on
Vivado.
\end{itemize}

\subsection*{Used by}

\code{razboj::sim::run} and \code{razboj::sim::run\_list}, and so Razboj's
tests, \code{//gpu/razboj:trace} and \code{//gpu/razboj:demo}. The system in
\code{soc}, where the rasteriser is the host at corner 1,0 of the
lattice.

\subsection*{Limits}

\code{//docs:razboj} lists them. It writes one pixel per cycle at most. It
has no depth buffer, so primitives are drawn in list order. It has no
interpolation, so a triangle is one colour, and no texture. It clips
no geometry, only boxes, and the assembler \code{Op::encode} does that
on the host. It has one read out at a time.

The code adds these. It has no reset port. A list holds at most 65535
entries, since \code{left} is sixteen bits. \code{BASE}, \code{DL}
and \code{CTRL} are constants of its type, so a design that draws
into two buffers moves the second one's rows past the first's, as the
flagship's icosahedron does (issue 986), rather than moving
\code{BASE}. A colour is 24 bits and its alpha 8.
